\chapter{革命、立宪与帝国终结}
\label{ch:revolution-and-empire-end}
\begin{figure}[htbp]
  \centering
  \begin{tikzpicture}[
    x=1cm,y=1cm,
    date/.style={font=\scriptsize\bfseries, text=dynasty},
    event/.style={draw=timeline, fill=paper, rounded corners=2pt, align=center,
      inner sep=3pt, font=\scriptsize, text=ink},
    note/.style={font=\scriptsize\color{source}, align=center}
  ]
    \draw[timeline, line width=1.2pt] (0,0) -- (12,0);
    \foreach \x/\year in {0/1908,2.0/1909,4.0/1910,6.0/1911.5,8.0/1911.10,10.0/1912.1,12/1912.2}
      {\draw[timeline] (\x,.12) -- (\x,-.12); \node[date, below] at (\x,-.18) {\year};}
    \node[event, above] at (0,.42) {光绪、慈禧去世；\\摄政体制};
    \node[event, below] at (2.0,-.62) {咨议局开办；\\路权争论加剧};
    \node[event, above] at (4.0,.42) {资政院开院；\\立宪咨询程序};
    \node[event, below] at (6.0,-.62) {铁路国有化；\\四川保路危机};
    \node[event, above] at (8.0,.42) {武昌起义；\\各地行动不同步};
    \node[event, below] at (10.0,-.62) {南京临时政府\\成立};
    \node[event, above] at (12,.42) {退位诏书；\\帝制名义终结};
    \node[note] at (6,-1.35)
      {宣言、军事控制、财政交接、谈判和日常秩序不能由一条“革命时间”同时说明。};
  \end{tikzpicture}
  \caption{帝国终结的编年定位（1908--1912，简表）。图将立宪程序、铁路危机、武昌起义、各省政治转换、南京政府与退位安排分列，以免把1911--1912年写成一夜完成的全国转向。主要据《清宣统政纪》、内阁及部院档、资政院与地方咨议局材料、铁路文件、电报和中外报刊；刊登、发报、宣言、军事行动与实际交接的日期和范围须分别核验。}
  \label{fig:revolution-and-empire-end-timeline}
\end{figure}

\begin{figure}[htbp]
  \centering
  \begin{tikzpicture}[
    x=1cm,y=1cm,
    place/.style={draw=timeline, fill=paper, rounded corners=2pt, align=center,
      inner sep=3pt, font=\small},
    court/.style={place, draw=dynasty, fill=dynasty!13},
    province/.style={place, draw=timeline, fill=timeline!13},
    revolutionary/.style={place, draw=source, fill=source!10},
    note/.style={align=center, font=\scriptsize\color{source}}
  ]
    \fill[paper] (-.7,-3.9) rectangle (7.85,3.85);
    \draw[dynasty, line width=.8pt] (-.7,3.85) -- (7.85,3.85);
    \node[font=\bfseries\large, text=dynasty] at (3.55,3.42)
      {1909--1912：立宪机构、铁路危机与省级转换};

    \node[court] (beijing) at (3.65,2.55) {北京\\资政院、部院\\谈判与退位};
    \node[province] (sichuan) at (.15,1.25) {四川\\保路抗议与\\地方军政危机};
    \node[revolutionary] (wuchang) at (2.55,.55) {武昌\\1911年起义\\新军与社团};
    \node[province] (hunan) at (1.1,-.8) {湖南及中部各省\\军政、财政与\\独立过程各异};
    \node[province] (shanghai) at (5.35,.35) {上海、江浙\\报刊、商会、\\交通与政权转换};
    \node[revolutionary] (nanjing) at (4.65,-1.3) {南京\\临时政府\\1912年初};
    \node[province] (guangdong) at (1.15,-2.65) {两广、福建及沿海\\宣言、军队与\\通信网络};
    \node[province] (northeast) at (6.7,1.95) {东北与北方\\袁世凯军政、\\财政交通与谈判};

    \draw[dynasty, dashed] (beijing) -- node[above,sloped,note]
      {诏令、谈判与军政调度} (northeast);
    \draw[->, source, line width=1.15pt] (sichuan) -- node[above,sloped,note]
      {危机在交通、财政与地方组织中传导} (wuchang);
    \draw[->, timeline, line width=1.15pt] (wuchang) -- node[above,sloped,note]
      {消息、军队与省级行动并非同步} (shanghai);
    \draw[->, timeline, dashed] (wuchang) -- node[left,sloped,note]
      {中部省份的不同转换} (hunan);
    \draw[->, source, line width=1.15pt] (shanghai) -- node[right,sloped,note]
      {城市、港口与代表联络} (nanjing);
    \draw[timeline, dashed] (hunan) -- node[below,sloped,note]
      {通信、人员与军队流动（示意）} (guangdong);
    \draw[dynasty, dashed] (beijing) to[out=-105,in=75] node[right,note]
      {退位与政权交接并行} (nanjing);

    \node[draw=source, fill=paper, rounded corners=3pt, align=center,
      text=source, font=\small] at (6.1,-2.8)
      {各省“独立”、军政控制、财政与交通\\不是同一时点或同一范围的事实};
  \end{tikzpicture}
  \caption{从立宪机构到帝国终结的省级政治空间（1909--1912，示意）。图以北京的立宪与退位程序、四川保路危机、武昌起义、上海--南京的政治联络及若干省级军政和通信节点为线索，强调宣言、军事控制、财政交通和谈判的时间并不重合。依据《清宣统政纪》、内阁与部院档、资政院及各省咨议局材料、铁路借款和路权文件、地方军政档、电报、中文外文报刊、革命与立宪团体刊物，并参照外使报告与各省地方材料。路线和位置不按比例；图中的省级节点只表示材料可见的政治联系或事件，不表示全省已一致“独立”、任何政权完成控制，亦不能由退位诏书推定社会秩序和财政交接已经结束。}
  \label{fig:revolution-provincial-order}
\end{figure}


\section{1909—1910：公共政治先有了新的房间}

1909年各省咨议局开办，给商绅、士人和新式团体一个可署名发言、记入议案并上呈北京的制度场域。它并不把省内所有人同样带入政治：资格、财产、教育、居住地以及通往省城的交通，仍决定谁能坐进议场，谁只能从报纸、告示或熟人口中得知议题。1910年资政院开院后，地方请愿与中央咨询之间多了一条可见的文书路线。可是同一年在若干地方出现的灾荒、粮价与财政压力提醒人们，能够进入程序的意见并不等于所有家庭都有余力参与程序；购粮、借贷、欠赋和迁徙，先在地方社会中改变了人们面对国家的方式。

新设议场也没有取代旧有的办事网络。州县的书吏与差役仍经手征收和文书，商会、会馆、宗族与学堂仍会筹款、调解或传递消息；督抚、部院和税务机关则掌握让议案变成命令的关节。因此，同一句“地方意见”在省城席位、商人账簿和乡村摊派之间并非同一件事。铁路尤其使这层差异难以遮掩：尚未铺成的路线已经对应已缴与待缴的股款、借款利息、征地、捐输和未来税源。中央所见的是可以整合的干线、抵押与借款，地方公司和股东所问的却是旧股怎样折算、负担是否仍留在本省、收益由谁处分。

咨议局章程、选举名册和会议记录可以说明制度何时开办、哪些人取得资格、哪些议案被提出；借款合同、公司账簿和税务文书则能追到产权与筹款如何写入纸面。它们都不足以把登记在册者化作全省民意，更不能从一项命令推出县乡的执行。商会档案、地方报刊、电报和地方志各自保存筹款者、城市读者、传讯者或后出编纂者的视角。正因材料的角度并不相同，立宪程序与社会困境应当并读，而不能各自写成一条不受对方牵动的时间线。

\section{1911：铁路、武昌与省城的不同节奏}

1911年铁路国有化并以外债筹资的决定，把国家希望整合干线、信用与借款的目标，撞上各省商办铁路、股东和地方精英对产权与税源的要求。\eventanchor{EQ-1911-RAILWAY-NATIONALIZATION}{铁路国有化（1911）}不是一个自动引发革命的按钮。在四川，保路运动、赵尔丰的处置和省城的消息网络使原可在请愿、会议与账目之间周旋的矛盾转成更尖锐的政治危机。对持股者，折价关系到已投入的钱；对曾为筑路捐输或受摊派者，问题可能是下一笔负担；对官府，则是维持秩序、上解财源和调兵成本无法轻易分开。地方志、公司账簿、督抚奏折和报纸各能说明一种位置，却不能被拼接成同一群体的意愿，或化约为一条单线因果。

十月武昌起义后，\eventanchor{EQ-1911-WUCHANG-UPRISING}{新军与革命团体控制武汉部分地区}。武汉首先面对的是兵员、枪械、渡江交通、粮饷和城内秩序；北京所能发出的调兵与请饷命令，则必须穿过电报、铁路、可用的官员和仍肯服从的部队。此前为四川危机而调兵改变了湖北一带的防务条件，却不能据此替每一名士兵指定政治选择，更不能把武昌归结为四川的单一后果。军事行动、财政可用性与消息传递在同一秋天相互牵动，缩短的是清廷从容处置危机的时间。

北京的命令、武汉的军队、南京的建政讨论与各省的独立声明没有在同一天发生，也不具有同一效力。宣布独立、守住省城、支付军饷和继续征收税粮是四件不同的事。省城改换旗号可以改变政治名义，却不能自行让县署、仓储、厘金局、邮电和团练按新的节奏运作；旧有的印信、档册与征收网络既可能被接收，也可能被军队、县衙或地方团体分别掌握。清廷起用袁世凯处理北方军政与议和，既是争取把军队、北京安全和谈判门径重新接合的选择，也承认宫廷无法单凭旧有链条指挥新军、省级军政和地方财政。

\section{1912：诏书结束名义，交接开始新的争执}

1912年南京临时政府成立，南北谈判与各地财政、军队和行政名义的交接并行。新的共和国名义不能先于可用的财政体系而自行落地：盐税、关税、厘金和地方杂税由谁保管，铁路账目与旧债由谁承认，军饷先付给哪一支部队，都会反过来决定任命、公报和改旗能否持续。旧官是否留任、商人是否垫款、士兵是否领到饷项，也不是中央文件可以预先写定的结果。革命的转折因而不是从清廷到共和国的一次替换，而是在省级机关、军营、税局与市场中同时发生、又彼此牵制的多次交接。

库伦方面在1911年末宣布脱离清廷，并拥立哲布尊丹巴为政教领袖；北部边疆的政教与外交关系随之进入另一套紧迫的谈判与行动节奏。这一变化不能被北京或南京的日历吞没。清方档案、蒙古文盟旗和寺院材料、俄方外交记录所使用的政治称谓、历法和关切并不相同，公告也不能单独画出稳定的控制范围。青藏高原、新疆、东北和西南的交通、地方权威与跨境关系同样各有待按本地材料辨认的转换过程。“帝国终结”因此不是一份中央文件由内地匀速铺向边地，而是不同道路、城镇、宗教网络与文书语言中的不同时点。

退位诏书为皇帝、部院和旗务的旧名义规定了制度终点：\eventanchor{EQ-1912-QING-ABDICATION}{清帝退位（1912）}。它没有结清军队听谁指挥、债务由谁承认、铁路和盐税怎样维持，也没有替旗人家庭、旧官员、商会和新式学校安排新的生计。对许多县乡居民，最早可见的变化未必是国号，而可能是征粮者、驻军者、县署告示和市场秩序的改变；这些变化既不在所有地方同时到来，也未必带来同一种结果。帝国的结束因此不是一幅末代皇帝离座的图像，而是行政、军事、财政和社会关系在不同时点被迫重写的开端。

\begin{textwindow}[退位诏书不是帝国经验的总计]
《宣统帝退位诏书》称：“今全国人民心理，多倾向共和。”这句话出自1912年谈判与交接中的官方文本，其作用是为退位提供一种能够被承接的制度语言。它可以说明清廷如何表述终局，却不是一次覆盖全国的意见调查：省份的军政控制、边疆社群的政治抉择和家庭的生计转换，仍须回到各地的文书、账目、报刊与其他材料中逐项追问。
\end{textwindow}

\paragraph{材料位置。} 《宣统政纪》、内阁与军机档、退位诏书、各省公报和临时政府文书能够核对程序、命令与刊布日期；它们首先是从发布者一端留下的行政链。四川保路档案、商会账簿、税册、军队给饷记录、地方报刊和日记可把视线移向筹款、交易、征收与消息抵达的具体地点；英日使领馆报告又保存外交利益所塑造的观察。革命派、立宪派和官报的宣传材料最适于辨认公开主张与动员语言，不宜直接等同于成员数、全省实际控制或群众态度。清历、蒙古历与格里历的换算亦须逐项核定，尤其不能把电报、公告或报纸的日期误当作整个地区同时完成政治转向的日期。

\section{立宪的场所与铁路的账本}

议场里的桌椅、选举名册和呈递北京的请愿书，并没有把晚清变成一个人人都能进入的公共空间；它们却给原先分散在会馆、商会、学堂和衙门门房之间的交涉，添上了可署名、可登报、可转寄的形式。1909年各省咨议局开办，1910年资政院成立。官员、绅商、学生和报人由此能够把地方事务说成“公议”，但资格、财产、教育、性别与居住地仍规定了谁能坐在场内，谁只能在场外听闻。

这种新场所并未替换旧有的治理网络。州县书吏继续经手赋役与文书，宗族、会馆和商会仍在筹款、调解和治安中保有作用；同一句“地方意见”，在省城议席、商人账册和乡村的差役摊派之间，可以指向很不相同的处境。铁路争论之所以很快压过程序本身，正在于它把这些层次连到了一起：一条尚未铺成的线路，已经牵动股份、借款、捐输、税源、征地和一省对自身资源的想象。

1911年的干线国有与外债筹资方案，不宜只读作中央和地方之间的一次口号对撞。中央希望把分散的线路、抵押和借款纳入可调度的财源；地方公司与股东则要追问已缴股款怎样折算、未来收益由谁处分、为路筹集的负担是否仍会留在本省。争执跨出会议厅，是因为它会落到商号的现金流、家庭的摊派和地方官的征收能力上。四川的抗议与赵尔丰的强制处置使危机骤然升级，却不能把铁路单独列为帝国终局的原因：灾荒、物价、军队调动、地方组织和革命网络同时改变了人们还能作出的选择。

电报和报刊缩短了消息抵达省城的时间，却没有制造同步的行动。发报、刊登、转抄、口传到沿线市镇和乡村，是一串不同速度的过程；一份城市报纸能记录某个社团的主张，不能替全省居民投票。立宪制度提供了争论的语言，也让中央债务、省级财政和地方要求无法彼此兼容的时刻更清楚地显露出来。

\subsection{材料与争议}

章程、选举名册、会议记录和内阁档能够追到制度何时开办、何种议案被提出；借款合同和公司账簿能够追到债务与股份如何被写入纸面。它们都不能自行证明一项命令在县乡怎样执行，更不能把登记在册的股东或议员直接化作“全省民意”。商会档案、地方报刊、电报和地方志保存的是不同位置的观察：前者常写筹款与交涉，后者更容易保留省城的声音，地方材料又会受编纂时间和政治立场影响。

因此，保路者的目标、参与范围和消息传播速度只能逐地检验。官报所称的秩序、抗议者宣称的响应和后出的革命叙事，都须同命令的下达、军队的移动、赋税能否收上来以及县城的交接分别对读。材料最清楚地显示的是制度动作和若干地点的冲突；它留下的空白，恰是普通人如何得知消息、是否参加以及付出何种代价最难整齐统计的部分。

\section{读者事件簇：省议场与中央咨询院}

\eventanchor{EQ-1909-PROVINCIAL-ASSEMBLIES}{各省咨议局开办（1909）}后，省城里原本在商会、学堂和官署间往返的人，有了按章程集会、记名发言和上呈意见的地方。它并非地方议会意义上的主权机关，议员资格也把多数妇女、贫民与远离省城的人留在门外；但对取得席位的商绅、士人和新式团体而言，地方财政、教育、铁路和官员作为开始可以被放进同一套公开程序中争论。

程序带来的并不是一条自下而上的命令线。议场若要把意见变成事情，仍须经过督抚、部院、税务机关与地方经手人；而州县征收、宗族调解、会馆筹款和商会交涉也不会因为章程生效而停止。选举和会议记录能看见谁被准许入场、谁提出何种主张，却看不见沉默者是否同意。正是在这层不完整的代表性上，“地方公议”获得了力量，也留下了它最容易被质疑的裂缝。

\eventanchor{EQ-1910-ZIZHENG-YUAN}{资政院开院（1910）}把咨询程序延伸到北京，也把省城的不满送进了一间距离税源、股款和粮价现场更远的屋子。代表、请愿书、报纸和电报在省议场与中央之间往返，能把一个问题变成全国可见的议题，却无法将各省的铁路股权、财政缺口和灾荒压力压缩成同一项意见。咨询之名给了交涉以门径，同时也标出了门槛：谁能递呈、谁能列席、何种议案能被采纳，仍由帝国的行政结构限定。

次年的铁路国有争议因而成为一次尖锐检验。它不是立宪机构“失败”之后才突然出现的外部事件，而是产权与财政冲突进入新程序、又从程序中溢出的过程。可续读\eventref{EQ-1911-RAILWAY-NATIONALIZATION}{铁路国有化}：当中央要以外债重整路权时，议席上的论辩会同公司账目、地方税收和街头动员一道承受后果。

\section{读者事件簇：路权争论先于国有化}

\eventanchor{EQ-1909-RAILWAY-RECOVERY-MOVEMENT}{路权与铁路国有化争论（1909）}已经使商办公司、地方股东和官府围绕产权与债务拉锯。对公司账房而言，线路是已收与待收的股款、尚未偿清的借款和可能的收益；对经手地方财政者而言，它又关系到为路筹款的名目能否收取、能否转作别用。铁路尚未把车开到许多地方，关于它的钱已经进入市镇与乡村的生活。

这也解释了“路权”为何不只是投资者的词汇。征地会改变沿线住户的安排，通往省城的想象会影响城镇商人，摊派或捐输则会使并不持股的人同样感到压力。不同人未必因这些处境而采取同一政治立场；一些人要求保留商办权，一些人忧虑债务与税负，还有一些人更在意地方资源是否会被调往外省。1909年的争论说明，次年出现的立宪议场可以收纳意见，却不足以消化已经嵌入财政和日常负担的分歧。

\section{读者事件簇：四川保路的省城压力}

四川保路运动在1911年升级，成都的请愿、商会与报刊活动、官府文书和赵尔丰的处置在同一座省城相遇，再沿电报线、商路和人际网络向外传去。“保路”把股东、商人、士绅、学生以及各地被卷入筹路负担的人暂时放进同一面旗帜下，却没有抹去他们对产权、税负、地方声望和政治出路的不同盘算。政府的强制也不能直接推知全省已被控制；它只能显示协商在关键地点让位于命令与拘捕。

成都的冲突之所以牵动全局，不在于它自动引发了别处的革命，而在于原本可在会议、文书与请愿之间周旋的压力失去了一层缓冲。清廷随后从湖北调兵入川，调动本身又改变了武汉一带新军的防务与人员配置。军队并不是一个无差别的工具：部队离开、命令到达、饷项能否跟上，以及士兵如何理解局势，都是不同的问题。四川危机与武昌起义之间存在这条可追踪的资源和时序联系，却不应被压缩成一根单线因果。

\section{读者事件簇：省份宣布独立的不同含义}

\eventanchor{EQ-1911-PROVINCIAL-INDEPENDENCE}{多省宣布脱离清廷（1911）}后，一纸独立公报与实际的军队听命、税粮征收、邮电运行和县乡交接常常不在同一天完成。省城改换旗号，可以切断旧朝在该处的名义；要让薪饷、粮食、巡防和文书继续流动，却仍要依靠旧官署的人员、地方商人的垫款和居民对风险的判断。每一份声明都让帝国网络再裂开一道口子，但它不是当地秩序已经重建的收据。

对地方官、商会、士兵和普通住户来说，“独立”的社会后果也不相同。有人面对的是效忠对象的转换，有人面对的是货物能否过关、店铺能否开门、家人能否避开征调。公报和回忆录特别容易留下宣告的声音；县乡的执行、拒绝、拖延和妥协则更零散。读这些省份变化时，应把政治名义、军事控制和日常生计分开，才不会把革命后的不整齐误作缺少历史意义。

\section{读者事件簇：袁世凯回到北方军政}

\eventanchor{EQ-1911-YUAN-SHIKAI-RETURN}{起用袁世凯（1911）}，是清廷试图把北方军队、谈判和财政重新接合起来的选择。朝廷需要的不只是一个官员的名望，而是一套可以传达命令、筹措军饷、同各方周旋的军政关系；袁世凯的回归因而改变了指挥链的重心。它没有把新军、省级军政与革命组织重新放回宫廷的单一节奏，也没有消除各地对饷项、归属和安全的现实顾虑。

议和由此不应写成单纯的礼仪退让。对北方军政集团、清廷和南方代表而言，谈判同时关涉兵权如何安置、债务由谁承接、税源何时能重新流动。新军在这场转折中既是被调遣的武装力量，也是由籍贯、军饷、训练和所在城市联系起来的具体人群；把他们笼统称作某一方的“工具”，会遮蔽命令失效时最先落到士兵和市民身上的风险。

\section{读者事件簇：南京临时政府的并行开端}

\eventanchor{EQ-1912-REPUBLIC-FOUNDING}{南京临时政府成立（1912）}时，南北谈判尚在推进，盐税、铁路、军饷和旧官署的交接已经不能等待一个抽象的“统一”完成。共和国的成立确立了新的中央名义，却没有使各地军队、边疆与城市在同一刻转换。铁路账目和税收的归属尤其提醒人们：宣布新的政府可以先于接收一套可用的财政与交通系统。

这一并行的开端把革命的社会尺度留在了日常事务里。地方若无稳定的饷源，士兵和家属便要面对不确定；商人若不知关卡、税捐和货运由谁负责，交易便难以恢复；旧官署的人员也须在新旧命令之间选择服从、观望或离去。临时政府的文件、谈判记录与各地公报可以标出新的政治语言，却不足以把这些不同的行政钟表调成同一时间。

\subsection{灾荒、物价与政治动员：一九一〇年的地方钟表}
\begin{event}{\eventanchor{EQ-1910-FAMINE-AND-UNREST}{多地灾荒、物价与不安}}{宣统二年（一九一〇年）}{清与多地社会}{灾荒和物价压力的报告可核；地区差异、死亡与政治态度不能由汇总数字或官员呈报概括}
一九一〇年的灾荒与物价压力先出现在地方的粮价、赈务、迁徙和欠赋问题中，再进入省级与中央的政治信息。它们没有在每一地都导向同一种抗争或革命选择；报告者的地域、统计口径和求援目的会改变所见。把这些压力直接译作“革命必然性”，反而会抹去地方社会如何处理危机的时间。

与咨议局、铁路产权的争执并读，这条线索提醒读者，晚清政治并不只发生在议场和电报线上。粮价上升时，家庭先要调整购粮、借贷和迁徙；地方绅商与官府则面对赈款从何而来、欠赋如何处理、道路是否仍能运粮的问题。财政、粮食和救济的可达性会改变谁有余力请愿、行商、从军或只是等待消息。各类报刊、档案和地方材料的限度见\hyperref[app:sources]{来源索引}。
\end{event}

\begin{event}{\eventanchor{EQ-1911-SICHUAN-RAILWAY-PROTEST}{四川保路升级：一条铁路怎样失去协商的缓冲}}{宣统三年（一九一一年）}{清与四川保路行动者 · 成都、川汉铁路沿线与省城电报网络}{抗议升级与赵尔丰处置可核；参与规模、组织目标和各地暴力过程不能由官报或报刊单独确定}
四川保路运动升级时，成都的请愿与会商已经不只是抽象的“反对国有”。川汉铁路的股款、地方税源、商办公司的账簿和外债安排，让商人、绅士、学生、股东与官府在同一条尚未完成的线路上承受彼此冲突的风险。对一名持股者，折价关系到已经投入的钱；对为筹路而被摊派的人，问题也许是下一笔负担能否承受；对地方官，维持秩序与上交财源又难以彼此分开。电报和报刊把消息送往州县，却不能让乡村、沿线城镇和省城同时作出同一种选择。

赵尔丰的处置把原本可在会议、文书和请愿之间周旋的争执推向强制。清廷从湖北调兵入川应付危机，既显示中央仍试图调动新军，也使武汉一带的驻防与军中关系发生变化。不能由这一调动直接推定每位士兵的政治选择，或把武昌起义归结为四川的单一后果；但军队、铁路与财政在同一秋天互相牵动，确实缩小了朝廷可以从容处置危机的时间。

《宣统政纪》的相关条目、督抚奏折、地方档案、公司材料和报刊能分别固定命令、请愿、消息发布或某一地点的冲突；它们的观察者与政治目标不同，不能合成一份全省一致的民意或伤亡统计。保路的直接后果是四川的财政与治安问题迅速政治化；中期影响则是武昌起义后的各省交接失去可供中央调度的兵力、信誉和时间。把四川看作帝国终局中的一条资源线，才不会把它降格为“革命爆发”前的一段背景：铁路既是未来的承诺，也是债务、税收与地方协商已经无法同时承受的现实。
\end{event}

\section{从武昌到退位：帝国在不同地方结束}

武昌起义使武汉三镇首先成为一处必须立刻处理兵员、枪械、渡江交通和城内秩序的战场。新军中的起事者和革命团体取得了部分城市空间，却还要面对清军的反攻、粮饷的筹集和消息向外传递的速度。朝廷所能看见的是一串战报、请饷和调兵的文书；城中与周边居民所面对的，则是军队驻扎、市场是否开张、家人能否通行以及哪一种命令会在次日仍然有效。这场起事因而不是一个抽象的“革命爆发”，而是把原先已因新军组织、革命网络和铁路风潮而绷紧的地方关系推入公开的军事竞争。

消息沿电报线、报纸和人际网络离开武汉，各省的回应却没有依同一张时刻表发生。有的省城先出现军队与官员的重新结盟，有的地方由商绅、咨议局人士和新式团体参与商议名义，有的地区则在省城宣布之后仍须逐县处理厘金、田赋、仓储和团练。所谓“独立”首先是一项政治宣言；它能否变成行政事实，要看谁能守住省城和交通线，谁能发给士兵军饷，谁能使税款不在征解途中断裂。旧督抚系统留下的印信、档册和征收网络可以被接管，也可以被地方军队、县衙与商会各自截留。因而，省份脱离北京的次序不能直接当作地方社会已经换政权的次序。

清廷起用袁世凯处理北方军政和议和事务，是在中央已难以单凭宫廷命令调动可信军队的条件下作出的选择。北洋军的指挥、北京的安全和与南方交涉的门径随之被放到同一张谈判桌上，但它们并不天然指向同一个结果。袁世凯既须向清廷说明军政处置，又要同革命方面周旋；革命方面也不能只以宣言取代军队、港口和城市财政的安排。战事与议和互相塑形：战场上的推进会改变谈判筹码，谈判的风声也会影响官员、士兵和地方团体究竟继续观望还是转向。

南京临时政府成立时，南北谈判尚未结束。它提出新的法理中心，却不能因此抹平各省军政和财政的实际差异。铁路国有化与外债引起的冲突已经显示，铁路股权、借款责任和地方筹款并非单纯的技术问题；辛亥之际，盐税、关税、厘金以及地方所收各项税款由谁保管、用来支付何处军队、是否承认旧债，同样决定了新名义能在何处落地。省级机关的改名、旧官的去留和军队的改编常在同一座城市里并行发生，不能把一纸任命或一份公报当成已经完成的交接。

北部边疆的时间也不服从北京与南京之间的叙事节奏。库伦方面在1911年末宣布脱离清廷，并拥立哲布尊丹巴为政教领袖；这一行动发生在清廷的财政、军政和通信同时承压之际，也牵动蒙古王公、宗教网络和跨境关系。公告能够表明一种政治主张，却不能单独画出稳定的控制范围，更不能把各地的支持写成整齐一致。青藏高原、新疆以及东北、西南的政局、交通与地方权威，各有尚待按本地材料厘清的转换过程。帝国的结束在边地不是从一份中央文件向外铺开的直线，而是多种旧关系和新主张在不同道路、城镇与文书语言中重新碰撞。

1912年的退位诏书为皇帝、部院和旗务的旧名义划出了明确的制度终点。南北谈判、财政军政压力和皇室待遇的安排共同促成这一结果；诏书所写的优待可以确认中央层面的承诺，却不能替代对各地兑现情形的追查。战争、债务、地方军队和身份供养并未随诏书消失。旗人家庭要面对生计与旧有供养的变动，旧官吏要决定是否留任，地方学校、商会和新式军队也要在新的制度名义下重新寻找资源与位置。对于许多县乡居民而言，最先可见的变化未必是国家称号，而可能是征粮者、驻军者、县署告示和市场秩序的改变；这些变化并不在所有地方同时到来，也不必然朝同一方向发展。

所以，帝国终结不是末代皇帝离开权位的一瞬，而是一组行政、军事、财政和社会关系被拆开重组的开端。退位解决了谁在中央名义上继承国家的问题，却把军队的归属、税源的流向、地方债务与边地权威留给此后的共和国继续承受。

\subsection{材料与争议}

退位诏书、内阁档案、各省公报与临时政府文书能把程序、命令和自我表述放回日序，却只能首先证明这些文书被拟定、传递或刊布。革命宣传、私人日记、报刊与外交报告保存的是动员者或观察者各自所处的位置；它们有助于看见消息如何流动，不能自动量出一省的支持程度。地方档案、商会账簿、税册和军队给饷记录更接近日常交接，但也常只留下一段机构的视野。

因此，省份独立的日期、地方实际控制、军队人数和群众态度须逐项互校。尤其是关于边疆，清方、俄方和蒙古文材料使用的称谓与历法并不一致；关于城市与县乡，电报抵达、报刊刊登、官署改名和税粮实际改征也不是同一件事。材料所能支持的是不同地点的局部时序，而不是一幅全国同步转向的图景。

\subsection{交叉阅读窗口：少年中国的排比，和1911年的不同时钟}

梁启超1900年刊行的《少年中国说》写道：“少年智则国智，少年富则国富。”这是一篇流亡改革派的政治散文，整齐排比把国家的将来托付给“少年”，以夸张、召唤式的节奏制造共同体的期待；它并非1911年武昌、四川或北京的现场报道，更不能替所有学生、士兵和乡村居民规定政治意愿。把它同《宣统政纪》、内阁档、各省公报和临时政府文书并读，能钉住政策、退位与交接程序；四川保路档案、商会账簿、地方报刊、《申报》与《民立报》则显示消息、产权和动员如何在不同城市流转，英日使领馆报告提供另一种受外交利益限制的观察。

钱穆的提问在于立宪、内阁和省级机构为何未能把新旧权力的交接稳定为共同程序；吕思勉的提问则在于债务、盐税、铁路股权、教育与地方组织怎样决定哪些人有条件进入公共政治。报刊可证明某日消息被刊登和以何种语气出现，不能证明全省同时知晓；电报和独立宣言也不等于县乡已经易手。梁启超的散文让我们听到革命前公共语言的音量，而多源互校才允许我们辨认帝国是在何处、以何种行政和社会节奏结束。


\section{读者事件簇：命令、谈判与退位没有同一速度}

\begin{event}{\eventanchor{EQ-1911-WUCHANG-UPRISING-DOCUMENTS}{武昌起义后的军政文书：北京如何追赶武汉的变化}}{宣统三年（一九一一年）}{清与革命军 · 武汉、北京及铁路电报线路}{军队调遣、饷械筹措和战报处理可由文书链核对；文书不能证明地方服从、民众态度或单一因果}
武昌的枪声传到北京后，清廷面对的是一条已经被打断的军政链：武汉的新军与革命团体控制了部分地区，中央却必须经由电报、铁路、军饷和仍肯执行的官员把回应送回长江中游。命令的发出不等于命令能穿过拥挤的线路、抵达各支部队，更不等于省城和县城会把同一消息理解成同一件事。筹措饷械、调动军队和编写战报，是清廷试图重新取得时间的行动；它也暴露出旧有中枢无法单独协调新军、地方财政和省级政治的限度。

《宣统政纪》宣统三年相关条、军机处档、战事方略及地方奏报可核对调遣与行政动作。它们大多站在发布命令或呈报局势的一端，不能据以断定武汉之外的实际控制、军队人数或普通居民的立场。此处形成的文书压力直接促成南北谈判的必要性：军事命令不能立刻恢复一致，政治交接便开始在不同城市同时展开。
\end{event}

\begin{event}{\eventanchor{EQ-1912-REPUBLIC-FOUNDING-DOCUMENTS}{南京临时政府成立：谈判与地方交接并行}}{宣统四年（一九一二年）}{中华民国临时政府与清 · 南京、北京及各省行政网络}{上谕、奏折和部议可证谈判及行政流转；不能由此推断政策有效、地方服从或所有行动者具有同一目标}
南京临时政府成立时，新的政治名义并没有让旧的税源、军队和官署自动换手。南京要把任命、财政和对外声明送往各省，北京则仍经由内阁、部院和袁世凯所能调动的军政网络处理谈判；同一时期，省城、县署和军营各自计算该使用何种印信、向谁请饷、如何避免眼前冲突。起义后的文书链在这里没有消失，而是分成相互竞争又必须接触的两套行政线路。临时政府的成立使“共和”成为可操作的谈判一方，也让交接的实际成本落到具体机构和地方关系上。

《宣统政纪》宣统四年相关条、宫中朱批奏折、内阁大库题本与《清史稿》相关志传保存了中央层面的程序。它们不足以证明每一省的改旗、军队归属或民间态度，档案的沉默不能由整齐的革命叙事填补。南京与北京并行处理的直接后果，是退位诏书必须同时考虑皇室、官署和军政接收；它也预示帝国终结后争夺的仍是人、财和命令的去向。
\end{event}

\begin{event}{\eventanchor{EQ-1912-QING-ABDICATION-DOCUMENTS}{退位诏书：制度终点与未结清的行政关系}}{宣统四年（一九一二年）}{清、袁世凯与革命派 · 北京宫廷、部院及各省交接线路}{退位诏书及相关上谕、奏折、部议可核；不能由其推断各地立即服从、群体态度或交接已经完成}
退位诏书把清帝国的最高名义置于终点，却没有把北京到各省的军队、盐税、债务和旗务一并结算。袁世凯、清廷与革命派在谈判中需要让一份文本既能说明皇室位置，也能给部院、旧官员和新政府留出继续办事的接口；对握有枪械、账册或地方关系的人来说，诏书只是必须据以重新判断听命对象的一项条件。它解决了谁以何种语言宣布终结的问题，却不能替不同地区安排同一天的政治经验。南京临时政府成立后并行的两套文书链，在这里被要求接成一套新的、仍不稳定的行政秩序。

《宣统政纪》宣统四年相关条、宫中朱批奏折、内阁大库题本与《清史稿》相关志传可核对退位程序与中央机构的表述。文书能够证明制度语言和可追查的行政动作，不能证明所有省份、军队、边疆社群或家庭已经完成转向。退位的直接效果是终止清朝皇统的正式名义；中期影响则是军事指挥、财政承认和地方治理继续在新的政治框架中争执，这也是“帝国终结”不能被缩成一纸诏书的原因。
\end{event}
\section{低密度区域事件簇}
\begin{event}{\eventref{EQ-1661-KANGXI-ACCESSION}{康熙继位}}{顺治十八年}{北京宫廷与文书网络}{继位可核；辅政分工不确定}
顺治帝去世后，八岁的玄烨继位，索尼、苏克萨哈、遏必隆和鳌拜以辅政大臣的身份进入朝政的枢纽。皇位的交接首先在北京的诏令、礼仪和官署中取得了可辨认的形式；但一个未成年的皇帝不能亲自把决定送往各处，诏令、旗务和南方战事所需的调发，便要经过辅政集团、部院与仍在运转的地方文书路线。继承因而不只是宫廷中的名分转换，也是哪些人能确认命令、调集资源并使其继续向外传递的问题。

实录、起居记录和遗诏档能够把继位及四名辅政者置于这一时间序列，也保留了朝廷希望公开的安排；它们并不能打开私下协商的全部过程，更不能把后来辅政者之间的权力消长预先写成既定结局。北京确立的集体入口须在八旗高层、部院程序和外省军政的实际往来中才会显出分量；皇权成长与旗人权力平衡的争执，也由此获得了日后展开的制度场所。
\end{event}
\begin{event}{\eventanchor{EQ-1683-TAIWAN-SURRENDER}{郑克塽投降}}{康熙二十二年}{澎湖、台湾海峡和府城}{投降可核；接收速度不确定}
澎湖战事之后，郑克塽向清廷投降，东宁政权的中枢由此结束。海峡上的胜负改变了府城可以同谁谈判、港口由谁接收的条件：郑氏一方原来依赖的海防与外援预期受到削弱，清廷则取得派员、安置官兵并讨论设治的政治入口。施琅及清廷所见的，首先是投降文书、人员处置和海防安排；这些安排必须跨过海峡，落实到驻军、港口和地方中介，才可能成为岛上的日常治理。

因此，接收不能被叙述成一夜完成的“归入”。府城的交接、军民安置、移民社会中的土地关系，以及不同原住民社群与沿海居民所面对的变化，并不共用一个速度，也未必留下同样密度的文字记录。实录与施琅奏折可见朝廷怎样界定投降和接收，后出的府志有助于追索设治后的地方框架；它们不足以代言全岛居民的态度，更不能把随后迁界调整、土地重组和贸易秩序的漫长过程压缩进投降当天。
\end{event}
\begin{event}{\eventanchor{EQ-1759-KASHGAR-YARKAND-CAMPAIGN}{清军入南疆}}{乾隆二十四年}{喀什噶尔、叶尔羌和军台}{进军可核；地方选择有争议}
乾隆二十四年，清军进入喀什噶尔、叶尔羌等南路绿洲，和卓政权的军事政治中心随之失去维系空间。这场进军并非只是在地图上越过天山。军队要沿军台和道路前行，补给要在遥远的转运线与绿洲之间接续；城镇中的地方力量、商路和供给关系，都会影响一支军队能否停驻和再向前推进。官修方略所称的胜利，因而首先标示清廷在战事中的推进与奖惩叙述，而不是对每一处绿洲都已稳定服从的直接证明。

战役打开的是长期治理的入口。驻军如何维持，税收如何征集，商路如何重新接上军台与绿洲，均要在战后逐处处理，并把成本持续送回中央与地方财政的网络。实录、军机处档、方略和地方奏报可以互相校对清军的调动与行政设想，却分别带着中央、军功与呈报的视角；它们不能据此精确量化伤亡，也不足以替代当地社群对战争、供给和新秩序所作的不同选择。天山南路纳入帝国的行政过程，正是在这些未被一场战役结清的关系中延长。
\end{event}
\begin{event}{\eventref{EQ-1842-NANJING-TREATY}{南京条约}}{道光二十二年}{南京江面、香港和五口}{签约可核；执行不能由文本推出}
长江航运受威胁、清军失利以后，清英代表在南京江面签订《南京条约》。文本把割让香港、五口通商和赔款等事项写成两国交涉的结果；它也把此前的战事压力转译为港口、关税、译介与财政上的持续事务。对清廷而言，签字不是战争后果的终点：条约要经过批准与换约，赔款须由财政渠道筹措，五个口岸的规定还要由地方官、港口人员、商人和译员在各自的工作中接住。

条约文本和照会能够说明双方在何种文字中提出权利与义务，实录和《筹办夷务始末》可追查清廷处理这一交涉的文书线索；它们不能从纸面推出每一港口何时、怎样运行，更不能证明居民或商人如何接受这些改变。香港的割让与五口开埠也不应被写成同一种地方经验：江面上的签约、跨海的殖民接管和各口岸的行政调整，分别依赖不同的机构、交通和财政条件。战争的结局由此没有消失，而是转化为长期而不均衡的执行压力。
\end{event}
\begin{event}{\eventref{EQ-1864-TIANJING-FALL}{天京陷落}}{同治三年}{天京、长江和江南粮运}{陷城可核；死伤不可确定}
长期围困使城内的粮食处境与太平天国的内部裂解日益尖锐；同治三年，湘军攻入天京，太平天国的核心政权遂告覆亡。攻城一刻能够改变的是城防和政权中枢的归属，并不能立刻修复江南在战争中断裂的田地、人口流动和征收关系。围城军队所依赖的饷项与转运、天京与长江航路之间的联结，以及城陷后人员的去留，都使战争的后果延伸到城外的州县、粮运和地方财政。

清廷实录、军机处档和地方奏报能够核对陷城及其后的处置语言，却也常以军功、收复和呈报的尺度整理这一结局。它们不能让战报中的伤亡、战果或战后秩序自行成为可接受的总数。对城内幸存者、逃离者和周边社会而言，恢复田亩、安置与重开征收是另一些缓慢而不整齐的过程；对承担军费的地方军事力量而言，战后如何保有筹饷与办事的地位又形成新的压力。天京陷落结束了一个政权中心，却没有结束江南社会为战争付出的代价。
\end{event}
\begin{event}{\eventanchor{EQ-1900-SOUTHEAST-MUTUAL-PROTECTION}{东南互保}}{光绪二十六年}{长江口、口岸和电报线}{互保可核；范围和动机不确定}
义和团战争期间，北京的战争命令与东南督抚面对的口岸、商贸和治安压力相撞。刘坤一、张之洞等督抚在电报往来和对外交涉中寻求避免战火南扩的办法，广东、长江流域及若干通商口岸的官员、商人和外方领事也各自掌握消息、航运与财政上的条件。“东南互保”所指的并不是一座独立政府的建立，而是在中央命令难以同速覆盖各地时，一组彼此联系却并不完全相同的地方应对。

朱批奏折、题本和实录可以追出清廷命令、督抚呈报与行政分歧的文书轨迹；它们最能说明谁向谁报告、谁以何种理由提出处置，并不能直接证明各省的范围整齐一致，更不能证明地方秩序已经稳定。电报线使消息和承诺可以迅速移动，却不能替代军队、港口管理、税收和地方协商所需的实际资源。互保由此暴露的，是中央—地方命令链出现了分叉；这种分叉没有取消清廷的名义，却使晚清危机中的区域选择和财政军事依赖更加清楚。
\end{event}
